\section{A Non-Linear Interlude: Counting Color Transitions}\label{sec:markov-interlude}

Before the main theorems, a short example worth walking through.
Section~\ref{sec:mergeable-sketches} gave the classical recipe for a tree
of summaries: encode shards, merge bounded states, and read out an
aggregate query. Section~\ref{sec:framework} generalized that recipe to a
learned summarizer under three local laws. The smallest task that exposes
why the learned generalization is useful is one whose oracle is a natural
aggregate but cannot be computed by an unordered register-wise merge unless
the state carries boundary information. This interlude walks through that
task and motivates why $g$ has to be a richer operator than a set-like
sketch merge.

\subsection{A Synthetic Document with a Non-Linear Oracle}\label{sec:color-sequences}

A document is a sequence of $128$ tokens. Each token is generated by a
two-step process: first, a hidden Markov chain draws one of $K$ colors;
then the color emits an observed token from a palette of $4$ synonyms.
Color~A might emit \emph{mist}, \emph{lake}, \emph{reed}; color~B emits
\emph{rust}, \emph{ember}, \emph{brick} (Figure~\ref{fig:markov-tree-walkthrough}).
The palettes are disjoint --- no token belongs to two colors --- but the
learner does not know which tokens share a color. It must discover the
groupings from data.

The oracle counts \emph{topic changes}: positions where the hidden color
switches from one token to the next. A typical document has about $5$
changes; Table~\ref{tab:dgp} gives the full data-generating specification.
This oracle is an integer aggregate just like a classical sketch query,
but it is \emph{not} a sum of per-token or per-leaf features. A topic
change at the boundary between two leaves is invisible to either leaf in
isolation: the left leaf sees only the token preceding the boundary, the
right leaf sees only the token following it, and neither can tell whether
those two tokens share a color without comparing them. The mergeable-sketch
recipe built on register-wise max or coordinate-wise add would drop the
boundary signal entirely.

\begin{figure}[t]
    \centering
    \includegraphics[width=\linewidth]{markov_changepoint_combined.pdf}
    \caption{\textbf{Counting topic changes: document and merge tree.}
    A $12$-token document with $4$ colors (A--D), each emitting synonym tokens (e.g., \emph{mist, lake, reed} for~A). The learner sees only the word row.
    Each leaf summary records (internal count, first color, last color). The merge adds children's counts plus $\One\{\text{left-last} \neq \text{right-first}\}$, detecting boundary changes invisible to either child.
    Naive per-leaf sum: $0+1+1+1=3$; tree: $5$ (correct).}
    \label{fig:markov-tree-walkthrough}
\end{figure}

\subsection{What Goes Wrong with an Unordered Merge}

Suppose we attempt to apply the HyperLogLog recipe directly: encode each
leaf into a register-like summary, merge pairwise with a commutative
operator, read out the count. For this to succeed the encode step would
have to reduce each leaf to a feature set that, under commutative merge,
still distinguishes between (i) a document whose two halves happen to end
and begin with the same color and (ii) a document whose two halves end
and begin with different colors. No register-wise max, additive count,
or bitmap-OR merge can distinguish these two cases: both produce the
same leaf-level aggregate. The only distinguishing information is
\emph{ordered and positional} --- which color is at the right edge of the
left leaf, which color is at the left edge of the right leaf --- and
a commutative merge by definition cannot preserve it.

The hand-designed sketch shown in Figure~\ref{fig:markov-tree-walkthrough}
works around this by explicitly carrying two pieces of boundary metadata
(first color, last color) in addition to the internal transition count.
The merge is then
\[
\text{merge}(L, R) \;=\; \Big(L.\text{count} + R.\text{count} + \One\{L.\text{last} \neq R.\text{first}\},\; L.\text{first},\; R.\text{last}\Big),
\]
which is associative but \emph{not} commutative. This works, but it
required a human to recognize that the oracle depends on ordered
boundaries and to hand-engineer a state that preserves them. For the
political-science application of Section~\ref{sec:manifesto-llm}, the
analogous engineering for RILE scoring would have to carry forward every
span-level cue the rubric touches --- a specification no one has written
down. We instead let $g$ \emph{learn} what boundary information to carry.

The learned $g$ we use throughout the rest of the paper is a Fourier neural operator---a position-aware architecture whose frequency-domain mixing naturally carries boundary context across the full sequence window. That architecture supplies representational capacity, not certification by itself. The certification target is still the local-law subspace: C1 asks each leaf to keep the internal count and boundary colors it needs, and C3 asks the merge to preserve the ordered boundary correction. Section~\ref{sec:fno-primer} introduces the realizability layer before Section~\ref{sec:markov-walkthrough} shows that a tree of FNO leaves solves this task and Section~\ref{sec:hll-parity} shows the same framework recovering the classical HyperLogLog sketch when the oracle happens to be a register-friendly query.
